\textbf{pTau217}: This is a blood measure that can reflect both amyloid and
tau-related changes in the brain. These changes can occur years or decades
before cognitive changes and alone do not determine if you have Alzheimer’s
disease. An increase in the blood test result levels suggests that microscopic
Alzheimer’s related changes may be present in the brain. This is not a
diagnostic test but in patients with memory changes this test can suggest that
memory problems may be related to Alzheimer’s changes in the brain.  In those
without memory problems, we are still studying if an abnormal test result can be
used to predict long term risk. False positive and false negative test results
occur.

Not Elevated: A value that is considered \textbf{not elevated} suggests against
the presence of amyloid plaques in the brain

Indeterminant: A value that is considered \textbf{indeterminant} suggests that
there is uncertainty regarding the presence of amyloid plaques in the brain.

Elevated: A value that is \textbf{elevated} suggests a higher risk for having
amyloid plaques in the brain.
